\chapter*{Introduction générale}
\addcontentsline{toc}{chapter}{Introduction générale}

L'essor des architectures conteneurisées a profondément transformé la manière dont les entreprises conçoivent, déploient et exploitent leurs applications. Kubernetes s'est imposé comme le standard de facto de l'orchestration de conteneurs, mais son adoption reste freinée par la complexité technique qu'il impose aux équipes de développement : gestion des espaces de noms, configuration de la mise à l'échelle, sécurisation des flux réseau entre applications, supervision de la charge. Parallèlement, le paradigme serverless, initialement porté par les offres propriétaires des grands fournisseurs de cloud public, a démontré l'intérêt d'une facturation alignée sur l'usage réel des ressources et d'une mise à l'échelle automatique pouvant aller jusqu'à l'arrêt complet d'une application en l'absence de trafic.

C'est dans ce contexte que s'inscrit ce Projet de Fin d'Études, réalisé au sein de NextStep IT, intégrateur spécialisé dans les solutions d'infrastructure et les services cloud. L'entreprise, dont l'offre Infrastructure as a Service met aujourd'hui à disposition de ses clients des ressources virtualisées brutes, souhaite enrichir cette offre par une couche d'abstraction supplémentaire : une plateforme PaaS (Platform as a Service) serverless multi-tenant, bâtie sur un cluster Kubernetes propre à l'entreprise, permettant à ses clients de déployer et d'exploiter des applications conteneurisées sans en gérer directement l'infrastructure sous-jacente.

L'objectif de ce mémoire est de présenter la démarche de conception et de réalisation de cette plateforme, depuis la mise en place du socle d'infrastructure jusqu'à la livraison d'un produit complet, intégrant un backend applicatif sécurisé, deux interfaces web, une chaîne d'intégration continue, et des mécanismes de supervision et de facturation à l'usage. Le projet s'appuie sur un ensemble de technologies open source reconnues : Kubernetes pour l'orchestration, Knative Serving et Knative Eventing pour les capacités serverless et événementielles, Apache Kafka opéré par l'opérateur Strimzi pour la messagerie asynchrone, Spring Boot et le client Fabric8 pour le backend et l'orchestration programmatique du cluster, React pour les interfaces utilisateur, Keycloak pour l'authentification, PostgreSQL pour la persistance, Jenkins et Kaniko pour l'intégration continue, ainsi que Prometheus et Grafana pour la supervision.

La conduite du projet a été organisée selon la méthodologie agile Scrum, retenue pour sa capacité à accompagner un développement itératif où chaque brique technique (infrastructure, puis serverless, puis messagerie événementielle, puis sécurité, puis fonctionnalités métier, puis interfaces graphiques, puis industrialisation) a pu être validée avant que la suivante ne s'y appuie. Cette organisation en releases et en sprints structure directement le plan de ce mémoire.

Ce mémoire est organisé comme suit. Le premier chapitre présente le contexte du projet, l'organisme d'accueil, la problématique traitée, une étude des solutions existantes sur le marché, ainsi que la méthodologie de conduite de projet retenue. Le deuxième chapitre détaille l'analyse des besoins fonctionnels et non fonctionnels de la plateforme, la modélisation globale du système (diagramme de cas d'utilisation, diagramme de classes, architecture logique et physique) et le pilotage du projet avec Scrum. Les chapitres suivants présentent, release par release et sprint par sprint, la réalisation progressive de la plateforme, depuis la mise en place de l'infrastructure Cloud Native jusqu'à la livraison des interfaces utilisateur et à l'industrialisation des chaînes de déploiement continu. Le dernier chapitre présente la validation de la plateforme, au travers des tests réalisés et d'un audit de sécurité en deux temps, ainsi qu'une synthèse honnête des limites du projet et de ses perspectives d'amélioration. Une conclusion générale clôt ce mémoire en revenant sur les apports du projet au regard de la problématique posée.
